\PassOptionsToPackage{unicode}{hyperref}
\PassOptionsToPackage{hyphens}{url}
\PassOptionsToPackage{dvipsnames,svgnames,x11names}{xcolor}
\documentclass[
  letterpaper,
  DIV=11,
  numbers=noendperiod]{scrartcl}
\usepackage{xcolor}
\usepackage{amsmath,amssymb}
\setcounter{secnumdepth}{-\maxdimen} % remove section numbering
\usepackage{iftex}
\ifPDFTeX
  \usepackage[T1]{fontenc}
  \usepackage[utf8]{inputenc}
  \usepackage{textcomp} % provide euro and other symbols
\else % if luatex or xetex
  \usepackage{unicode-math} % this also loads fontspec
  \defaultfontfeatures{Scale=MatchLowercase}
  \defaultfontfeatures[\rmfamily]{Ligatures=TeX,Scale=1}
\fi
\usepackage{lmodern}
\ifPDFTeX\else
\fi
\IfFileExists{upquote.sty}{\usepackage{upquote}}{}
\IfFileExists{microtype.sty}{% use microtype if available
  \usepackage[]{microtype}
  \UseMicrotypeSet[protrusion]{basicmath} % disable protrusion for tt fonts
}{}
\makeatletter
\@ifundefined{KOMAClassName}{% if non-KOMA class
  \IfFileExists{parskip.sty}{%
    \usepackage{parskip}
  }{% else
    \setlength{\parindent}{0pt}
    \setlength{\parskip}{6pt plus 2pt minus 1pt}}
}{% if KOMA class
  \KOMAoptions{parskip=half}}
\makeatother
\makeatletter
\ifx\paragraph\undefined\else
  \let\oldparagraph\paragraph
  \renewcommand{\paragraph}{
    \@ifstar
      \xxxParagraphStar
      \xxxParagraphNoStar
  }
  \newcommand{\xxxParagraphStar}[1]{\oldparagraph*{#1}\mbox{}}
  \newcommand{\xxxParagraphNoStar}[1]{\oldparagraph{#1}\mbox{}}
\fi
\ifx\subparagraph\undefined\else
  \let\oldsubparagraph\subparagraph
  \renewcommand{\subparagraph}{
    \@ifstar
      \xxxSubParagraphStar
      \xxxSubParagraphNoStar
  }
  \newcommand{\xxxSubParagraphStar}[1]{\oldsubparagraph*{#1}\mbox{}}
  \newcommand{\xxxSubParagraphNoStar}[1]{\oldsubparagraph{#1}\mbox{}}
\fi
\makeatother

\usepackage{longtable,booktabs,array}
\newcounter{none} % for unnumbered tables
\usepackage{calc} % for calculating minipage widths
\usepackage{etoolbox}
\makeatletter
\patchcmd\longtable{\par}{\if@noskipsec\mbox{}\fi\par}{}{}
\makeatother
\IfFileExists{footnotehyper.sty}{\usepackage{footnotehyper}}{\usepackage{footnote}}
\makesavenoteenv{longtable}
\usepackage{graphicx}
\makeatletter
\newsavebox\pandoc@box
\newcommand*\pandocbounded[1]{% scales image to fit in text height/width
  \sbox\pandoc@box{#1}%
  \Gscale@div\@tempa{\textheight}{\dimexpr\ht\pandoc@box+\dp\pandoc@box\relax}%
  \Gscale@div\@tempb{\linewidth}{\wd\pandoc@box}%
  \ifdim\@tempb\p@<\@tempa\p@\let\@tempa\@tempb\fi% select the smaller of both
  \ifdim\@tempa\p@<\p@\scalebox{\@tempa}{\usebox\pandoc@box}%
  \else\usebox{\pandoc@box}%
  \fi%
}
\def\fps@figure{htbp}
\makeatother

\setlength{\emergencystretch}{3em} % prevent overfull lines

\providecommand{\tightlist}{%
  \setlength{\itemsep}{0pt}\setlength{\parskip}{0pt}}

\KOMAoption{captions}{tableheading}
\makeatletter
\@ifpackageloaded{caption}{}{\usepackage{caption}}
\AtBeginDocument{%
\ifdefined\contentsname
  \renewcommand*\contentsname{Table of contents}
\else
  \newcommand\contentsname{Table of contents}
\fi
\ifdefined\listfigurename
  \renewcommand*\listfigurename{List of Figures}
\else
  \newcommand\listfigurename{List of Figures}
\fi
\ifdefined\listtablename
  \renewcommand*\listtablename{List of Tables}
\else
  \newcommand\listtablename{List of Tables}
\fi
\ifdefined\figurename
  \renewcommand*\figurename{Figure}
\else
  \newcommand\figurename{Figure}
\fi
\ifdefined\tablename
  \renewcommand*\tablename{Table}
\else
  \newcommand\tablename{Table}
\fi
}
\@ifpackageloaded{float}{}{\usepackage{float}}
\floatstyle{ruled}
\@ifundefined{c@chapter}{\newfloat{codelisting}{h}{lop}}{\newfloat{codelisting}{h}{lop}[chapter]}
\floatname{codelisting}{Listing}
\newcommand*\listoflistings{\listof{codelisting}{List of Listings}}
\makeatother
\makeatletter
\makeatother
\makeatletter
\@ifpackageloaded{caption}{}{\usepackage{caption}}
\@ifpackageloaded{subcaption}{}{\usepackage{subcaption}}
\makeatother
\usepackage{bookmark}
\IfFileExists{xurl.sty}{\usepackage{xurl}}{} % add URL line breaks if available
\urlstyle{same}
\hypersetup{
  pdftitle={Problem Set 8: DNA in a nucleus},
  pdfauthor={Jun Allard},
  colorlinks=true,
  linkcolor={blue},
  filecolor={Maroon},
  citecolor={Blue},
  urlcolor={Blue},
  pdfcreator={LaTeX via pandoc}}

\ifXeTeX
\special{pdf:trailerid [ <25c8a98c3832990558cd1e5e1fe289c7> <25c8a98c3832990558cd1e5e1fe289c7> ]}
\fi
\ifPDFTeX
\pdftrailerid{}
\pdftrailer{/ID [ <25c8a98c3832990558cd1e5e1fe289c7> <25c8a98c3832990558cd1e5e1fe289c7> ]}
\fi
\ifLuaTeX
\pdfvariable trailerid {[ <25c8a98c3832990558cd1e5e1fe289c7> <25c8a98c3832990558cd1e5e1fe289c7> ]}
\fi

\title{Problem Set 8: DNA in a nucleus}
\usepackage{etoolbox}
\makeatletter
\providecommand{\subtitle}[1]{% add subtitle to \maketitle
  \apptocmd{\@title}{\par {\large #1 \par}}{}{}
}
\makeatother
\subtitle{Physics 230A/146A --- Biological Physics}
\author{Jun Allard}
\date{}
\begin{document}
\maketitle


DNA in eukaryotic cells is packed by complex machinery inside the nucleus.
Here, will use the ideal chain model to explore the need for this machinery by studying what would happen if the DNA were not carefully packed, i.e., by assuming it is just a freely-jointed chain randomly packed in a spherical nucleus.

\begin{enumerate}
\def\labelenumi{\roman{enumi}.}
\item
  Find the following quantities. Cite your sources.

  \begin{itemize}
  \tightlist
  \item
    The total length of DNA of a typical human chromosome, \(l\),
  \item
    The size (roughly the diameter) of a typical eukaryotic nucleus, \(D\),
  \item
    The width of double-stranded DNA, \(w\).
  \item
    The persistence length of double-stranded DNA, \(l_p\).
  \end{itemize}
\end{enumerate}

We will use these quantities to estimate how big a chromosome is, first from the point-of-view of occupation fraction, then from the point-of-view of polymer physics.

\begin{enumerate}
\def\labelenumi{\roman{enumi}.}
\setcounter{enumi}{1}
\tightlist
\item
  Assume DNA is a long cylinder of width \(w\) and length \(l\).
  What is the total volume of DNA in each cell?
  You may assume a single chromosome.
\item
  Assuming the nucleus is spherical, what fraction of the volume of the nucleus is occupied by DNA?
\end{enumerate}

For the remainder of the problem, you may use order-of-magnitude precision.

The end-to-end distance \(r_{ee}\) of a flexible polymer satisfies the probability distribution

\[
p(r_{ee}) = 4\pi r_{ee}^2 \left( \frac{3}{2\pi l\,l_p}\right)^{\frac{3}{2}} \exp\left(\frac{-3r_{ee}^2}{2l\,l_p}\right).
\]

If a system explores macrostates parametrized by \(r\), and in the absence of confinement or external forces has probability distribution \(p(r)\), then the entropy of macrostate \(r\) is

\[
S(r) = k_B \ln \left( p(r)\right) + S_0
\]

where \(S_0\) is a constant that is independent of the system's state.
(Eventually, only derivatives of \(S(r)\) matter for forces and pressures, so these should be independent of \(S_0\).)

If a polymer is confined to a sphere of diameter \(D\), this would mean the maximum distance between any two nodes, \(r_\mathrm{max}\), is less than \(D\).
Therefore, to be rigorous, computing the force a polymer exerts by this confinement would require the probability distribution \(p_\mathrm{max}(r_\mathrm{max})\).
Unfortunately we do not have an expression for \(p_\mathrm{max}(r_\mathrm{max})\).
Let us instead make an approximation that there exists a characteristic polymer ``size'' \(r\) such that \(r \sim r_{ee} \sim r_\mathrm{max}\).

\begin{enumerate}
\def\labelenumi{\roman{enumi}.}
\setcounter{enumi}{3}
\tightlist
\item
  Write an expression for the free energy \(G = -T S\) as a function of polymer size \(r\).
\item
  Pressure is force per unit area.
  Find an expression for the outward pressure exerted by the polymer if it is confined to a sphere of radius \(r=R\).
\item
  Use your findings to estimate the pressure exerted by DNA inside the nucleus, if it were randomly packed.
  Report your value in units of atmospheres (1 atmosphere \(\approx 10^5\) Pascals).
\item
  We wish to answer the question: is the pressure you found in (vi) sufficient to rupture the nucleus?
  What number should this pressure be compared with?
  Look up this number and make this comparison.
\end{enumerate}




\end{document}
